\section{Successor receiving identities on the original programme}

The preceding edition's MRI1--11 remain preserved. This section supplies
their more precise successors. The scalar statements below retain the
stipulated actual quartet, the full multiplicity $m_0$, $k=4\ell+1$, and
$q=[1+k(m_0-1)](k+1)^2$. The correlated native-row statements retain their
own simple-quartet degree $a=k+4$, dimension $q_a=(a+1)^2$, conductor
$\Delta=16a-48$, and the two original source orders. These are the
original domains of SR and CRS; no degree or source-order substitution
is made between them.

\subsection{The scalar, source order and complete action}
Write $d=4m_0-2$, $c_\zeta=2\gamma_E-\log(2\pi)$ and retain the physical
equilibrium constant $C_\Gamma$ and the complete-source pressure integral
$I_h$. The full SR1--30 proof gives, on its explicit finite guards,
\[
 \delta_{k,1}^{\rm ar}=-dC_\Gamma q+
 \frac{\log2}{\pi}I_hk^2+
 \frac{C_\Gamma q}{2(\log q+c_\zeta)}+R_{h,k}.
 \tag{ARU1}
\]
Here $R_{h,k}$ is exactly the remainder in SR29, with the SR30 bound
$o_h(q/\log q)$. Inserting ARU1 in the unchanged MRI3 identity proves
\[
\begin{split}
 J_k={}&\mathcal B_k^{(1)}+\mathcal W_k^{\rm mono}
 -\mathcal P_{k,-}-\mathcal P_{k,+}-dC_\Gamma q\\
 &+\frac{\log2}{\pi}I_hk^2+
 \frac{C_\Gamma q}{2(\log q+c_\zeta)}+R_{h,k}.
\end{split}\tag{ARU2}
\]
Subtracting the original source-order-one and source-order-$k$ identities
for this same $J_k$ gives the exact map
\[
 \delta_{k,k}^{\rm ar}=\delta_{k,1}^{\rm ar}
                   +\mathcal B_k^{(1)}-\mathcal B_k^{(k)}.
 \tag{ARU3}
\]
All source masses in the two baselines stay in this difference. On the
simple stratum define, as before,
$\mathfrak R_{h,k}=J_k-\delta_{k,1}^{\rm ar}-C_Bq^2-8q\log k$.
The previously proved estimate is $\mathfrak R_{h,k}=O_h(q)$.
Using this definition and ARU1, without absorbing either remainder,
proves
\[
\begin{split}
 J_k={}&C_Bq^2+8q\log k-2C_\Gamma q+
 \frac{\log2}{\pi}I_hk^2+
 \frac{C_\Gamma q}{2(\log q+c_\zeta)}\\
 &+\mathfrak R_{h,k}+R_{h,k}.
\end{split}\tag{ARU4}
\]
Consequently ARU1's improved arithmetic remainder does not improve
the separate analytic remainder $\mathfrak R_{h,k}$ by substitution.
Their exact relation is ARU4.

The full TS1--41 proof establishes $0<C_\Gamma<2/3$ and
$I_h>1681/(50\pi)$ from the actual theta source with its complete
denominator and outer mass. Its simple-stratum coefficient satisfies
\[
 d_0(h)=-2C_\Gamma+(\log2/\pi)I_h
       >7370383/7260000>1.
 \tag{ARU5}
\]
The finite guard TS40 therefore proves $\delta_{k,1}^{\rm ar}>q$.
For every fixed $m_0\ge2$, divide ARU1 by its literal $q$.
Then $k^2/q\to0$, $(\log q+c_\zeta)^{-1}\to0$ and
$R_{h,k}/q\to0$, proving the negative limit
$-(4m_0-2)C_\Gamma$. These are signs of the specified scalar on the
specified multiplicity strata. The exact projected receiver appears
below and retains its additional variables.

\subsection{The two path remainders and every marked frame}
Use the full two-parameter path
$F(t,r)=\mathcal B_k[\sigma e^{tH_C+rH_O}]$ of SR41--46.
Taylor's integral identity in $t$ and the fundamental theorem in $r$
give respectively
\[
\begin{split}
 F(1,1)-F(0,1)&=F_t(0,1)+
                  \int_0^1(1-t)F_{tt}(t,1)\,dt,\\
 F_t(0,1)-F_t(0,0)&=\int_0^1F_{tr}(0,r)\,dr.
\end{split}\tag{ARU6}
\]
Adding proves the original mixed and nonlinear receiver
\[
 \mathcal C_k^{\rm mixed}+\mathcal N_k^{\rm nonlinear}
 =F(1,1)-F(0,1)-F_t(0,0),\qquad
 |\mathcal C_k^{\rm mixed}+\mathcal N_k^{\rm nonlinear}|
 \le E_{\rm cen}+E_{\rm lin}=o_h(q/\log q).
 \tag{ARU7}
\]
The finite expressions for both errors, rather than just their order,
are SR5 and SR43. The matrices inside $F_{tr}$ retain the full
mixed trace SR44; the matrices inside $F_{tt}$ retain the full
projection variance SR45. Thus this bound applies to their displayed
sum with its original signs.

For an invertible original marked frame $\mathcal J$,
$\det(\mathcal J^*G_N[\nu]\mathcal J)=
 |\det\mathcal J|^2\det G_N[\nu]$. Apply this twice, with the same
$\mathcal J$ and the original arithmetic and Gamma measures. Division
proves
\[
 \log\frac{\det(\mathcal J^*G_N[m_k]\mathcal J)}
                 {\det(\mathcal J^*G_N[\sigma]\mathcal J)}
 =\log\frac{\det G_N[m_k]}{\det G_N[\sigma]}.
 \tag{ARU8}
\]
The same equality holds in each signed endpoint occurrence; summing
transports ARU1, with its unchanged error, into the marked relative
determinant. MAR supplies the exact regularized frame when the naive
frame degenerates. Its explicit inclusion and the attained metric are
used before ARU8; a singular determinant is never cancelled.
GMF supplies the growing-flag high/low comparison and retains its
terminal--low Schur block. FCD's fixed-degree expansion is used only
on its stated fixed-degree domain.

\subsection{Corrected native rows and their backward cancellation}
In the original CRS domain write $\gamma_r$ for the attained row
logarithm. The complete CRP and CRS derivations prove
\[
\begin{split}
 c_0&=\frac\Delta2\{1+\log(q_a/\Delta)\},\\
 c_r&=\Delta f(r/q_a)-\frac12
       \{(r+\Delta)\log(1+\Delta/r)-\Delta\}\quad(r\ge1),\\
 \sum_r|\gamma_r-c_r|&\le E_c=O_{\rm actual}(q_a).
\end{split}\tag{ARU9}
\]
The summation range and every finite guard are those of CRS17,
including the extension through $q_a+32a$. For weights $w_r$ equal
to the original endpoint weights $1,2,\ldots,2,1$, the triangle
inequality gives $|\sum w_r(\gamma_r-c_r)|\le2E_c$.
CRS23 retains the complete equilibrium integral for each prefix and
tail. At $R_*=\lfloor a^{3/2}/\log a\rfloor$, the deterministic
corrections to those two integrals are respectively
\[
 -64q_a\log a+128q_a\log\log a+O(q_a),\qquad
 -64q_a\log a-128q_a\log\log a+O(q_a).
 \tag{ARU10}
\]
Their sum is the existing $-128q_a\log a+O(q_a)$. Adding the
original lower-root $+128q_a\log a$ cancels it. This proves that the
more precise prefix/tail calculation introduces no extra
$q_a\log a$ term into ARU4. Retaining the old uncorrected
$\Delta f$ while merely reducing its error would omit exactly
the correction calculated in ARU10.

For the actual row projection $p_r$ and
$\kappa_r=-\log p_r$, retain $\kappa_r=+\infty$ when $p_r=0$.
CRP33 gives the two clipped expressions
\[
 \widetilde S_K=\sum_r w_r\min(c_r^+,\kappa_r),\qquad
 \widetilde S_R=\sum_r w_r(c_r^+-\kappa_r)_+,
 \quad |S_X-\widetilde S_X|
 \le2E_c+(\log2)\sum_r w_r\quad(X=K,R).
 \tag{ARU11}
\]
The pointwise equality
$\min(u,v)+(u-v)_+=u$ for $u,v\ge0$, also valid for $v=+\infty$
under the stated convention, proves
$\widetilde S_K+\widetilde S_R=\sum_rw_rc_r^+$.
The actual determinant identity remains $S_K+S_R=\mathcal T$.
The eight original signed flags retain their two ordered chains:
CRP28 gives the radial Lipschitz constant $4$ and row clipping
allowance $8\log2$; CRP29 gives $4\log2$ for each complementary
CPQ pair. Their original projection spectra are still present
in these formulas. This is the exact receiving map for their
next evaluation; no allocation is supplied by rank alone.
For the independent proper source, CRS27--30 likewise update
the displacement on its original four intervals, while retaining
its standard-window determinant.

\subsection{Complete retained class and projected arithmetic current}
NCB and NCRP prove the original-minor invertible coordinate completion
$F_I=[F_0,F_C,L_a]$ and the exact decomposition
$v_\lambda=F_0u_0+F_C\eta+L_az$.
For its full coordinate vector $w=(u_0,\eta,z)$, define
$H=F_I^*G_NF_I$ and $A_I=F_I^{-1}AF_I$. Matrix multiplication gives
\[
 w^*(A_I^*H+HA_I-aH)w
 =v_\lambda^*(A^*G_N+G_NA-aG_N)v_\lambda.
 \tag{ARU12}
\]
This identity includes all nine blocks of $H$. Expanding either side
after inserting the three summands proves equality term by term,
including both directed cross terms. It propagates the completed
class to the original arithmetic receiver without altering the native
section or its quotient maps.

For the literal observation projection $P$ and original boundary rows,
NCE proves
\[
 \ell_N=b_N^*G_N/\omega_N,\quad \xi_N=-G_Nb_{N+1},\qquad
 U=\frac{\ell_NPv_\lambda}{\ell_Nv_\lambda},\quad
 V=\frac{\xi_N^*Pv_\lambda}{\xi_N^*v_\lambda}.
 \tag{ARU13}
\]
Substituting the full three-component decomposition in each numerator
and denominator is exact because the two rows are linear. The
nonzero complementary components therefore remain in both ratios.
With the original current $c$ and $E=\|v_\lambda\|_{G_N}^2$, the
complete projected current is
\[
 \frac{\Phi_X}{2E}=\Re(cP_X),\quad
 (P_K,P_B,P_\times)=
 (\bar UV,(1-\bar U)(1-V),\bar U+V-2\bar UV),\quad
 P_K+P_B+P_\times=1.
 \tag{ARU14}
\]
Expanding the three products proves the last equality and hence
$\sum_X\Phi_X=2E\Re c$. This also proves the same-class cancellation
after the full-class repair, including both leakage directions.

JRE1--22 prove the exact joint image metric and its attained remainder.
Parity gives $F_{12}=0$ on this same quotient. Put
$p_j=|F_{j3}|^2/(EF_{jj})$, $\theta=K_{33}/E$ and retain the
original finite current allowance $E_{\rm cur}$. The resulting bound is
\[
\begin{gathered}
 d_*^2=\theta(1-\theta)(p_1^{-1}+p_2^{-1}),\qquad
 h_*={\theta(1-\theta)\sqrt{1-p_1-p_2}\over2\sqrt{p_1p_2}},\\
 (|\mathcal R_K|,|\mathcal R_B|,|\mathcal R_\times|)
 \le2EE_{\rm cur}
 (\theta d_*+h_*,(1-\theta)d_*+h_*,|1-2\theta|d_*+2h_*),
 \qquad \sum_X\mathcal R_X=0.
\end{gathered}\tag{ARU15}
\]
The componentwise bounds follow by inserting JRE21 in JRE17;
the cancellation follows from the imaginary part of ARU14.
The factors $p_1,p_2,\theta$ are original projection energies with
$p_1,p_2>0$, $p_1+p_2\le1$, $0\le\theta\le1$, proved in JRE20.
Their actual values and the actual directional spectra in ARU11
are the next analytic calculation. All earlier completed scalar,
period, marked-frame and finite-field calculations remain available
through the explicit maps above.
